\chapter{INTRODUCTION}
\section{General}
Nepal is located in a highly active seismic region where damaging earthquakes have repeatedly affected the built environment. The 1934 Bihar--Nepal earthquake and the 2015 Gorkha earthquake are among the events that exposed the vulnerability of buildings and reinforced the need for improved seismic assessment of important facilities \cite{gautam2016,adhikari2022}. Hospitals require particular attention because structural damage, excessive deformation, or loss of nonstructural functionality can reduce emergency medical capacity precisely when demand for those services is greatest \cite{guragain2004,subedi2022}.

Structural models commonly idealize the base of a building as fixed. That assumption is convenient, but it neglects the flexibility of the supporting soil and foundation. Soil--structure interaction (SSI) treats the superstructure, foundation, and soil as a coupled system and can alter the effective period, seismic force distribution, lateral deformation, settlement, and member demand \cite{wolf1985,poudel2024}. The influence of SSI becomes more important as soil and foundation flexibility increase.

This thesis examines a G+5 reinforced-concrete mid-size general hospital building in Nepal under four support conditions: fixed base, hard soil, medium soil, and soft soil. ETABS response-spectrum analysis provides the principal linear comparison. SAP2000 pushover results are used to interpret nonlinear static capacity, while nonlinear time-history results from twelve ground motions are used for a preliminary fragility assessment. The available results do not specify the analysed response direction, and one Barpark/Barpak-labelled record could not be independently verified; the fragility assessment is therefore treated as preliminary.
